\documentclass{article}
\usepackage[T1]{fontenc}
\usepackage{lmodern}
\usepackage{graphicx}
\usepackage{amsmath,amssymb}
\usepackage[a4paper,margin=2cm]{geometry}
\usepackage[absolute,overlay]{textpos}
\setlength{\TPHorizModule}{1mm}
\setlength{\TPVertModule}{1mm}
\usepackage[indonesian]{babel}
\usepackage{hyperref}
\setlength{\emergencystretch}{2em}
% unit: Halaman 52

\begin{document}

% === Halaman 52 (llm) ===

\section*{Keuntungan ECC}

\begin{itemize}
    \item Keuntungan yang sama dengan sistem kriptografi lain: \textit{confidentiality}, \textit{integrity}, \textit{authentication} and \textit{non-repudiation}, tetapi...
    \item Panjang kuncinya lebih pendek
    \begin{itemize}
        \item Mempercepat proses \textit{encryption}, \textit{decryption}, dan \textit{signature verification}
        \item Penghematan \textit{storage} dan \textit{bandwidth}
    \end{itemize}
\end{itemize}

\vspace{50mm}

\textcolor{red}{*) Sumber bahan: Debdeep Mukhopadhyay, Elliptic Curve Cryptography , Dept of Computer Sc and Engg IIT Madras}

\end{document}
